\documentclass[../main.tex]{subfiles}

\begin{document}

\subsection{Vector Iterations and the QR Algorithm}
\label{sub:eigen-algorithms}

\textBox{%
Section~\ref{sub:eigen-theory} characterized eigenvalues but did not compute them. Every polynomial is the characteristic polynomial of its companion matrix, so the Abel--Ruffini theorem implies that no formula built from finitely many arithmetic operations and root extractions gives the eigenvalues of a general matrix with $n\geq5$. Every eigenvalue algorithm is therefore iterative. We first study iterations that apply $\bA$ or $(\bA-\sigma\bI)^{-1}$ to a single vector. These iterations compute one eigenpair at a time and form the basis of the sparse methods in later sections. We then describe the QR algorithm, which computes all eigenvalues of a dense matrix by a sequence of orthogonal similarity transformations.
}

\subsubsection{Power Iteration and Shifted Inverse Iteration}

\addDef{Power Iteration}{%
Let $\bA$ be real symmetric. Given a unit vector $\bv^{(0)}$, power iteration constructs
\begin{align*}
    \bw^{(k)} = \bA \bv^{(k-1)}, \qquad \bv^{(k)} = \frac{\bw^{(k)}}{\|\bw^{(k)}\|_2}, \qquad \lambda^{(k)} = (\bv^{(k)})^T \bA \bv^{(k)}.
\end{align*}
Order the eigenvalues by magnitude. If $|\lambda_1|>|\lambda_2|\geq\cdots\geq|\lambda_n|$ and $\bq_1^T\bv^{(0)}\neq0$, the direction of the iterate converges linearly:
\begin{align*}
    \|\bv^{(k)} - (\pm \bq_1)\|_2 = O\left( \left| \frac{\lambda_2}{\lambda_1} \right|^k \right).
\end{align*}
For a general nonnormal matrix, an eigenvector basis may be poorly conditioned or may not exist, and this simple orthogonal expansion no longer justifies the claim.
}
\label{def:power-iteration}

\addDef{Shifted Inverse Iteration}{%
For a real symmetric matrix, an eigenvalue near a target $\sigma$ can be isolated by applying power iteration implicitly to $(\bA-\sigma\bI)^{-1}$:
\begin{align*}
    (\bA - \sigma \bI) \bw^{(k)} = \bv^{(k-1)}, \qquad \bv^{(k)} = \frac{\bw^{(k)}}{\|\bw^{(k)}\|_2}.
\end{align*}
The eigenvalues of $(\bA - \sigma \bI)^{-1}$ are $(\lambda_i - \sigma)^{-1}$. Provided the shift is not itself an eigenvalue and the starting vector contains the desired component, the mode closest to $\sigma$ dominates. The matrix inverse is never formed. Each step solves a shifted linear system instead.
}
\label{def:inverse-iteration}

\addDef{Rayleigh Quotient Iteration (RQI)}{%
For symmetric matrices, updating the shift dynamically to the current Rayleigh estimate $\sigma_k = (\bv^{(k)})^T \bA \bv^{(k)}$ defines Rayleigh Quotient Iteration:
\begin{align*}
    (\bA - \sigma_k \bI) \bw^{(k+1)} = \bv^{(k)}, \qquad \bv^{(k+1)} = \frac{\bw^{(k+1)}}{\|\bw^{(k+1)}\|_2}, \qquad \sigma_{k+1} = (\bv^{(k+1)})^T \bA \bv^{(k+1)}.
\end{align*}
}
\label{def:rqi}

\addThm{Local Cubic Convergence of Rayleigh Quotient Iteration}{%
Let $\bA$ be real symmetric, and let $\lambda_*$ be a simple eigenvalue with unit eigenvector $\bq_*$. If a unit RQI iterate $\bv^{(k)}$ is sufficiently close to $\bq_*$, then
\begin{align*}
    \sin\angle(\bv^{(k+1)},\bq_*)
    \leq C\sin^3\angle(\bv^{(k)},\bq_*).
\end{align*}
If $\lambda_*$ is an extremal eigenvalue, the Rayleigh shifts also satisfy
\begin{align*}
    |\sigma_{k+1}-\lambda_*|
    \leq C'|\sigma_k-\lambda_*|^3
\end{align*}
for iterates sufficiently close to $\bq_*$.
}
\label{thm:rqi-cubic}

\addPrf{%
Write the current unit vector as
\begin{align*}
    \bv=\cos\theta\,\bq_*+\sin\theta\,\bz,
    \qquad
    \bz\perp\bq_*,
    \qquad
    \|\bz\|_2=1.
\end{align*}
Its Rayleigh quotient $\sigma=\bv^T\bA\bv$ obeys
\begin{align*}
    \sigma
    &= (\cos\theta\,\bq_*+\sin\theta\,\bz)^T\bA(\cos\theta\,\bq_*+\sin\theta\,\bz) \\
    &= \cos^2\theta\,\lambda_* + 2\cos\theta\sin\theta\,\bq_*^T\bA\bz + \sin^2\theta\,\bz^T\bA\bz.
\end{align*}
Because $\bA\bq_* = \lambda_*\bq_*$ and $\bq_*^T\bz = 0$, the cross-term vanishes: $\bq_*^T\bA\bz = \lambda_*\bq_*^T\bz = 0$. Substituting $\cos^2\theta = 1-\sin^2\theta$,
\begin{align*}
    \sigma-\lambda_*
    = (1-\sin^2\theta)\lambda_* + \sin^2\theta\,\bz^T\bA\bz - \lambda_*
    = \sin^2\theta\,(\bz^T\bA\bz-\lambda_*),
\end{align*}
so $|\sigma-\lambda_*|\leq C_0\sin^2\theta$.

Let $\gamma=\min_{\lambda_j\neq\lambda_*}|\lambda_j-\lambda_*|>0$. Once $\theta$ is small enough, $|\lambda_j-\sigma|\geq\gamma/2$ for every other eigenvalue. Expanding the inverse-iteration vector $\bw=(\bA-\sigma\bI)^{-1}\bv$ in an orthonormal eigenbasis, its component along $\bq_*$ has magnitude
\begin{align*}
    \frac{|\cos\theta|}{|\lambda_*-\sigma|},
\end{align*}
whereas the norm of its orthogonal component is at most $2|\sin\theta|/\gamma$. Normalization does not change their ratio. Therefore,
\begin{align*}
    \tan\angle(\bv^+,\bq_*)
    \leq
    \frac{2|\lambda_*-\sigma|}{\gamma|\cos\theta|}|\sin\theta|
    \leq C_1|\sin\theta|^3.
\end{align*}
For small angles, $\sin\theta$ and $\tan\theta$ are comparable, which proves cubic convergence of the eigenvector direction.

If $\lambda_*$ is extremal, the spectral gap also gives two-sided local bounds $c\sin^2\theta\leq|\sigma-\lambda_*|\leq C\sin^2\theta$. Applying the Rayleigh-quotient identity to the next iterate and using the cubic angle bound yields
\begin{align*}
    |\sigma^+-\lambda_*|
    \leq C_2\sin^2\angle(\bv^+,\bq_*)
    \leq C_3\sin^6\theta
    \leq C_4|\sigma-\lambda_*|^3.
\end{align*}
}

\subsubsection{The Practical QR Algorithm}

\addDef{Basic QR Iteration}{%
Given $\bA_0 = \bA$, the unshifted QR iteration generates orthogonally similar matrices via
\begin{align*}
    \bA_k = \bQ_k \bR_k, \qquad \bA_{k+1} = \bR_k \bQ_k = \bQ_k^T \bA_k \bQ_k.
\end{align*}
All iterates share the same eigenvalues. Under suitable separation conditions the subdiagonal entries tend to zero and the iterates reveal a Schur form. Shifts and deflation are essential for a reliable practical method.
}
\label{def:qr-iteration}

\addLem{Invariance of Hessenberg Structure under QR Iteration}{%
Let $\bH \in \fR^{n \times n}$ be an upper Hessenberg matrix ($H_{ij} = 0$ for all $i > j + 1$). If $\bH = \bQ\bR$ is the QR factorization of $\bH$, then the reversed product $\bH_+ = \bR\bQ$ is also upper Hessenberg.
}
\label{lem:hessenberg-qr-invariance}

\addPrf{%
Because $\bH = \bQ\bR$ and $\bR$ is upper triangular and invertible (assuming $\bH$ is nonsingular, or extending by continuity), $\bQ = \bH \bR^{-1}$. Because the inverse of an upper triangular matrix is upper triangular, $\bR^{-1}$ is upper triangular. The product of an upper Hessenberg matrix $\bH$ and an upper triangular matrix $\bR^{-1}$ is upper Hessenberg. Therefore, $\bQ$ is upper Hessenberg.

Now consider $\bH_+ = \bR\bQ$. The matrix $\bR$ is upper triangular and $\bQ$ is upper Hessenberg. For entry $(i, j)$ with $i > j + 1$:
\begin{align*}
    (\bH_+)_{ij} = \sum_{k=1}^n R_{ik} Q_{kj}.
\end{align*}
Because $\bR$ is upper triangular, $R_{ik} = 0$ for all $k < i$. For the remaining terms $k \geq i$, we have $k \geq i > j + 1$, which implies $k > j + 1$. Because $\bQ$ is upper Hessenberg, $Q_{kj} = 0$ for all $k > j + 1$. Thus every term in the sum vanishes, proving $(\bH_+)_{ij} = 0$ for $i > j + 1$.
}

\addThm{Practical Shifted QR Framework}{%
A standard dense nonsymmetric eigenvalue computation has two phases:
\begin{enumerate}
    \item \textbf{Hessenberg reduction:}\enspace Compute $\bH_0=\bU_0^T\bA\bU_0$ with Householder reflectors. This one-time reduction costs $O(n^3)$ operations and preserves the eigenvalues.
    \item \textbf{Implicit shifted QR steps:}\enspace Apply shifts while preserving Hessenberg form. A schematic single-shift step is
    \begin{align*}
        \bH_k - \mu_k \bI = \bQ_k \bR_k, \qquad \bH_{k+1} = \bR_k \bQ_k + \mu_k \bI.
    \end{align*}
    Lem.~\ref{lem:hessenberg-qr-invariance} explains why each step costs $O(n^2)$ rather than $O(n^3)$. Real nonsymmetric codes use implicit single or double shifts so that complex conjugate pairs can remain in real arithmetic. A small subdiagonal entry permits deflation and splits the problem into smaller blocks.
\end{enumerate}
For a real symmetric matrix, Hessenberg form reduces to tridiagonal form. Wilkinson shifts from the trailing two by two block and implicit bulge chasing then give a particularly efficient symmetric algorithm.
}
\label{thm:practical-qr-algorithm}

\subsubsection{The Jacobi Eigenvalue Algorithm}

\textBox{%
For a symmetric matrix, an alternative to the QR algorithm drives the off-diagonal entries to zero directly. Each step applies an orthogonal similarity by a plane rotation chosen to annihilate one off-diagonal entry. The method, due to Jacobi in 1846, is slower than shifted QR for a general dense matrix. It remains useful because it can compute the small eigenvalues of graded positive definite matrices to high relative accuracy \citep{demmel1997}, and because rotations acting on disjoint pairs of indices can be applied in parallel.
}

\addDef{Givens Plane Rotation Matrix}{%
The matrix $\bG(p, q, \theta) \in \fR^{n \times n}$ is an identity matrix except for the four entries at the intersections of rows and columns $p$ and $q$:
\begin{align*}
    g_{pp} = c, \qquad g_{pq} = s, \qquad g_{qp} = -s, \qquad g_{qq} = c,
\end{align*}
where $c = \cos\theta$ and $s = \sin\theta$. The transformation $\bA' = \bG^T \bA \bG$ modifies only rows and columns $p$ and $q$ of $\bA$.
}
\label{def:givens-rotation}

\addThm{Derivation of the Optimal Rotation Angle}{%
To zero the transformed off-diagonal entry $a'_{pq} = 0$, we consider the $2 \times 2$ orthogonal similarity transformation on the submatrix at indices $\{p, q\}$:
\begin{align*}
    \begin{pmatrix} a'_{pp} & a'_{pq} \\ a'_{pq} & a'_{qq} \end{pmatrix}
    &=
    \begin{pmatrix} c & -s \\ s & c \end{pmatrix}
    \begin{pmatrix} a_{pp} & a_{pq} \\ a_{pq} & a_{qq} \end{pmatrix}
    \begin{pmatrix} c & s \\ -s & c \end{pmatrix} \\
    &=
    \begin{pmatrix}
        c^2 a_{pp} - 2cs a_{pq} + s^2 a_{qq} & (c^2 - s^2) a_{pq} + cs (a_{pp} - a_{qq}) \\
        (c^2 - s^2) a_{pq} + cs (a_{pp} - a_{qq}) & s^2 a_{pp} + 2cs a_{pq} + c^2 a_{qq}
    \end{pmatrix}.
\end{align*}
Setting $a'_{pq} = 0$ yields the trigonometric condition
\begin{align*}
    (c^2 - s^2) a_{pq} + cs (a_{pp} - a_{qq}) = 0
    \implies
    \frac{\cos(2\theta)}{\sin(2\theta)} = \frac{a_{qq} - a_{pp}}{2a_{pq}}.
\end{align*}
Defining $\tau = \cot(2\theta) = \frac{a_{qq} - a_{pp}}{2a_{pq}}$ and using $\cot(2\theta) = \frac{1 - \tan^2\theta}{2\tan\theta}$, setting $t = \tan\theta$ produces
\begin{align*}
    t^2 + 2\tau t - 1 = 0.
\end{align*}
To minimize the perturbation to $\bA$ and avoid cancellation, we select the root of smaller absolute value:
\begin{align*}
    t = \frac{\operatorname{sign}(\tau)}{|\tau| + \sqrt{1 + \tau^2}}, \qquad c = \frac{1}{\sqrt{1 + t^2}}, \qquad s = ct.
\end{align*}
}
\label{thm:jacobi-angle-derivation}

\addPrf{%
The quadratic formula yields roots $t = -\tau \pm \sqrt{1 + \tau^2}$. Multiplying the root of smaller magnitude by its algebraic conjugate yields
\begin{align*}
    t = \frac{(-\tau \pm \sqrt{1 + \tau^2})(\tau \pm \sqrt{1 + \tau^2})}{\tau \pm \sqrt{1 + \tau^2}} = \frac{1}{\tau + \operatorname{sign}(\tau)\sqrt{1 + \tau^2}} = \frac{\operatorname{sign}(\tau)}{|\tau| + \sqrt{1 + \tau^2}}.
\end{align*}
This formula satisfies $|t| \leq 1$, corresponding to a rotation angle $|\theta| \leq \pi/4$, and avoids subtracting nearly equal floating-point numbers.
}

\textBox{%
A rotation that annihilates $a_{pq}$ may make a previously annihilated entry nonzero again, so the method does not terminate after $n(n-1)/2$ steps. Convergence follows instead from a quantity that decreases at every step.
}

\addThm{Monotone Decrease of the Off-Diagonal Frobenius Norm}{%
Let the off-diagonal squared norm be defined by
\begin{align*}
    \operatorname{off}(\bA)^2 = \sum_{i \neq j} a_{ij}^2 = \|\bA\|_F^2 - \sum_{i=1}^n a_{ii}^2.
\end{align*}
Applying the Jacobi rotation $\bA' = \bG(p, q, \theta)^T \bA \bG(p, q, \theta)$ with the optimal angle strictly reduces the off-diagonal norm:
\begin{align*}
    \operatorname{off}(\bA')^2 = \operatorname{off}(\bA)^2 - 2a_{pq}^2.
\end{align*}
}
\label{thm:jacobi-convergence}

\addPrf{%
Because $\bG$ is orthogonal, the total Frobenius norm is preserved: $\|\bA'\|_F^2 = \|\bA\|_F^2$. The rotation affects only entries in rows and columns $p$ and $q$. For the $2 \times 2$ submatrix, the sum of squares is invariant under orthogonal similarity:
\begin{align*}
    (a'_{pp})^2 + (a'_{qq})^2 + 2(a'_{pq})^2 = a_{pp}^2 + a_{qq}^2 + 2a_{pq}^2.
\end{align*}
Because the rotation enforces $a'_{pq} = 0$, the sum of squared diagonal entries increases:
\begin{align*}
    (a'_{pp})^2 + (a'_{qq})^2 = a_{pp}^2 + a_{qq}^2 + 2a_{pq}^2.
\end{align*}
All other diagonal entries $a_{ii}$ ($i \neq p, q$) are untouched. Substituting this into $\operatorname{off}(\bA')^2 = \|\bA'\|_F^2 - \sum (a'_{ii})^2$ proves that $\operatorname{off}(\bA')^2 = \operatorname{off}(\bA)^2 - 2a_{pq}^2$.
}

\addDef{Cyclic Jacobi Algorithm}{%
Let $\bA^{(0)}=\bA$ and $\bV=\bI_n$. While $\operatorname{off}(\bA)>\texttt{tol}\,\|\bA^{(0)}\|_F$, sweep through the pairs $1\leq p<q\leq n$ in a fixed order. For each pair, compute $c$ and $s$ from Thm~\ref{thm:jacobi-angle-derivation}, replace $\bA$ by $\bG(p,q,\theta)^T\bA\bG(p,q,\theta)$, and replace $\bV$ by $\bV\bG(p,q,\theta)$. On termination, the diagonal of $\bA$ approximates the eigenvalues and the columns of $\bV$ the eigenvectors.
}
\label{def:cyclic-jacobi}

\textBox{%
Theorem~\ref{thm:jacobi-convergence} shows that $\operatorname{off}(\bA)^2$ decreases by $2a_{pq}^2$ at every rotation. Choosing the largest off-diagonal entry at each step gives linear convergence, as the exercises show. The cyclic order avoids that search, and once the off-diagonal entries are small compared with the gaps between the eigenvalues, cyclic Jacobi converges quadratically, measured sweep by sweep \citep{golubvanloan2013}. Each sweep costs $O(n^3)$ flops.
}

\subsubsection{Shifted Systems from One Hessenberg Reduction}

\textBox{%
Several calculations require the solution of $(s\bI-\bA)\bx=\bb$ for many complex shifts $s$ with the same $\bA$ and $\bb$. The vector $(s\bI-\bA)^{-1}\bb$ is the resolvent of $\bA$ applied to $\bb$. Resolvents appear in Laplace transforms of solutions of linear differential equations, in frequency responses, and in contour-integral formulas for matrix functions. Factoring $s\bI-\bA$ separately for each of $N$ shifts costs about $\tfrac23n^3$ flops per shift. The Hessenberg reduction of Thm~\ref{thm:practical-qr-algorithm} removes this repetition, because adding a multiple of the identity does not change Hessenberg structure.
}

\addThm{Shifted Solves with a Hessenberg Matrix}{%
Let $\bA=\bQ\bH\bQ^T$ with $\bQ$ orthogonal and $\bH$ upper Hessenberg, and let $s$ not be an eigenvalue of $\bA$. Then $\bx=(s\bI-\bA)^{-1}\bb$ is given by
\begin{align*}
    (s\bI-\bH)\by=\bQ^T\bb,
    \qquad
    \bx=\bQ\by.
\end{align*}
Gaussian elimination with partial pivoting solves the Hessenberg system in $O(n^2)$ operations. At each step, the only possible row interchange is between the pivot row and the row directly below it.
}
\label{thm:hessenberg-shifted-solve}

\addPrf{%
Since $\bQ\bQ^T=\bI$, we have $s\bI-\bA=\bQ(s\bI-\bH)\bQ^T$, so $\bx=\bQ(s\bI-\bH)^{-1}\bQ^T\bb$. The matrix $s\bI-\bH$ differs from $-\bH$ only on the diagonal, so it is upper Hessenberg. Suppose that after $k-1$ elimination steps the active submatrix is upper Hessenberg. Its first column then has a single nonzero entry below the diagonal, in row $k+1$. Partial pivoting compares only rows $k$ and $k+1$, computes one multiplier, and updates one row of length at most $n-k+1$. Rows $k+2,\dots,n$ are untouched, so the next active submatrix is again upper Hessenberg. Summing the work over $k=1,\dots,n-1$ gives $O(n^2)$ operations, and the triangular solve that follows costs another $O(n^2)$.
}

\textBox{%
Reducing $\bA$ to Hessenberg form with Householder reflectors costs about $\tfrac{10}{3}n^3$ flops when $\bQ$ is kept as a product of reflectors, and this reduction is done once \citep{golubvanloan2013}. Applying $\bQ^T$ and $\bQ$ to a vector costs $O(n^2)$ flops. For $N$ shifts, the total cost is therefore about $\tfrac{10}{3}n^3+O(Nn^2)$ flops, compared with about $\tfrac23Nn^3$ for separate factorizations. The reduction pays for itself after five shifts, and with one hundred shifts it reduces the cost by a factor of about twenty.
}

\addRemark{%
\textbf{(Contour integrals for the matrix exponential)}\enspace
The Cauchy integral formula for matrix functions gives
\begin{align*}
    e^{t\bA}\bb=\frac{1}{2\pi\mathrm{i}}\int_\Gamma e^{ts}(s\bI-\bA)^{-1}\bb\,ds,
\end{align*}
where $\Gamma$ is a contour that encloses the spectrum of $\bA$ \citep{higham2008}. When the spectrum lies in the left half-plane, as for Markov generators and discretized diffusion operators, $\Gamma$ can be chosen to wrap around the negative real axis so that $e^{ts}$ decays along it. After a parametrization of $\Gamma$, the trapezoid rule with $N$ nodes approximates the integral by a sum of $N$ resolvent applications, each computed by Thm~\ref{thm:hessenberg-shifted-solve}. For suitable Talbot-type contours the error decreases geometrically in $N$, and a few tens of nodes commonly reach double-precision accuracy \citep{weidemantrefethen2007}. The geometric convergence is the same mechanism as the spectral accuracy of the trapezoid rule for periodic analytic integrands in Section~\ref{subsec:quadrature}. For large sparse matrices, the Krylov approximations of Section~\ref{sub:arnoldi} are usually preferable, because they avoid the dense Hessenberg reduction.
}

\subsubsection{Exercises}

\addExcr{%
\begin{enumerate}
    \item \exCompute \textbf{(Observed cubic convergence of RQI)}\enspace Apply RQI to a symmetric matrix with distinct eigenvalues, starting from several unit vectors close to the same eigenvector.
    \begin{enumerate}
        \item Record the angle errors $\theta_k=\angle(\bv^{(k)},\bq_*)$ and the ratios $\sin\theta_{k+1}/\sin^3\theta_k$.
        \item Record the Rayleigh errors and verify the cubic relation in Thm~\ref{thm:rqi-cubic} before rounding error dominates.
        \item Repeat with two nearly repeated eigenvalues and explain how the spectral gap enters the constants in the proof.
    \end{enumerate}
    \item \exExplore \textbf{(Failure of unshifted QR on an exchange matrix)}\enspace Let $\bP = \begin{pmatrix} 0 & 1 \\ 1 & 0 \end{pmatrix}$.
    \begin{enumerate}
        \item Compute the QR factorization $\bP = \bQ_0 \bR_0$.
        \item Compute $\bP_1 = \bR_0 \bQ_0$ and show that $\bP_1 = \bP$.
        \item Repeat the step with a nonzero shift. Determine which shifts break the cycle, and explain how an eigenvalue shift from the trailing two by two block permits immediate deflation in this example.
    \end{enumerate}
    \item \exCompute \textbf{(Givens rotation Hessenberg QR)}\enspace Write a Python function that executes one step of shifted QR on an upper Hessenberg matrix $\bH \in \fR^{n \times n}$ using $n-1$ Givens rotations. Verify that your implementation requires $O(n^2)$ flops rather than $O(n^3)$.
    \item \exCompute Perform one Jacobi rotation by hand to zero out the off-diagonal entry of
    \begin{align*}
        \bA = \begin{pmatrix} 3 & 2 \\ 2 & 6 \end{pmatrix}.
    \end{align*}
    Explicitly compute $\tau, t, c, s$, and the resulting diagonal matrix $\bA'$.
    \item \exCompute Verify that the eigenvalues computed on the diagonal of $\bA'$ match the exact roots of $\det(\bA - \lambda\bI) = 0$.
    \item \exDerive Prove that for an $n \times n$ matrix, if the largest off-diagonal entry is selected at each step (Classical Jacobi), then $\operatorname{off}(\bA^{(k+1)})^2 \leq \left( 1 - \frac{2}{n(n-1)} \right) \operatorname{off}(\bA^{(k)})^2$.
    \item \exExplore \textbf{(Many shifted solves)}\enspace For a random $200\times200$ matrix, compare the time to solve $(s_k\bI-\bA)\bx_k=\bb$ for $k=1,\dots,N$ by separate LU factorizations and by one Hessenberg reduction followed by Thm~\ref{thm:hessenberg-shifted-solve}. Find the value of $N$ at which the two approaches cost the same, and compare it with the flop estimate in the text.
\end{enumerate}
}
\label{ex:eigen-algorithm-exercises}

\end{document}
